\expandafter\ifx\csname CAPETINCLUDEMODE\endcsname\relax
\documentclass[11pt,a4paper]{article}
% preamble.tex — packages et commandes communes a toutes les fiches CAPET SII IE
\input{glyphtounicode}
\pdfgentounicode=1
\usepackage{cmap}
\usepackage[T1]{fontenc}
\usepackage[utf8]{inputenc}
\usepackage[french]{babel}
% babel[french] active a \begin{document} des raccourcis de ponctuation (:;!?)
% qui cassent la syntaxe to[...] de circuitikz (erreur "+ or - expected") ;
% on les desactive, ce qui ne change que l'espacement fin avant ces signes.
\AtBeginDocument{\shorthandoff{:;!?}}
\usepackage[a4paper, top=1.6cm, bottom=1.6cm, left=1.4cm, right=1.4cm]{geometry}
\setlength{\headheight}{14.5pt}
\usepackage{amsmath, amssymb, amsfonts}
\usepackage{siunitx}
\sisetup{locale=FR, output-decimal-marker={,}, inter-unit-product=\ensuremath{{}\cdot{}},
         group-minimum-digits=4}
% Unites non fournies par siunitx, declarees ici une fois pour toutes (fiches + exercices).
% Les gardes \ifdefined evitent tout conflit avec une redeclaration locale restee dans un
% fichier : ne pas retirer les gardes, ne pas redeclarer ces unites dans les fichiers.
\DeclareSIUnit{\jour}{j}
\ifdefined\var\else\DeclareSIUnit{\var}{var}\fi   % puissance reactive
\ifdefined\vad\else\DeclareSIUnit{\vad}{VAD}\fi   % puissance deformante
\ifdefined\tr\else\DeclareSIUnit{\tr}{tr}\fi      % tour (tr/min)
\usepackage{array, booktabs, multirow, tabularx}
\usepackage{xcolor}
\usepackage[most, breakable]{tcolorbox}
\usepackage{multicol}
\usepackage{enumitem}
\usepackage{mathtools}
\usepackage{lmodern}
% Refonte visuelle 2026 (direction "moderne coloré") : Lato pour le texte,
% maths sans-serif assorties via newtxsf (charge apres amsmath/amssymb).
\usepackage[default]{lato}
\usepackage{newtxsf}
\usepackage{microtype}
\usepackage{fancyhdr}
\usepackage{tikz}
\usetikzlibrary{positioning,arrows.meta}
\usepackage{circuitikz}
\usepackage{lastpage}

% ── Couleurs par domaine (palette modernisée 2026) ────────────────────
\definecolor{colElec}{RGB}{29, 78, 216}        % Électrotechnique — bleu vif
\definecolor{colElecLight}{RGB}{219, 234, 254} % fond bleu clair
\definecolor{colAuto}{RGB}{21, 128, 61}        % Automatique — vert
\definecolor{colMeca}{RGB}{194, 65, 12}        % Mécanique — orange
\definecolor{colTherm}{RGB}{185, 28, 28}       % Thermique — rouge
\definecolor{colQT}{RGB}{232, 244, 233}        % (hérité, plus utilisé par questiontype)
\definecolor{colJury}{RGB}{255, 248, 225}      % fond jury (crème)
\definecolor{colJuryFrame}{RGB}{202, 138, 4}   % bordure/bandeau jury (ambre)
\definecolor{colGrayLight}{RGB}{245, 245, 247} % fond gris neutre
\definecolor{colSeparator}{RGB}{180, 180, 180} % gris séparateur
\definecolor{colExo}{RGB}{8, 145, 178}          % Banque d'exercices — cyan (distinct des fiches)
% ── Encre et gris de la refonte ───────────────────────────────────────
\definecolor{colInk}{RGB}{23, 32, 42}          % texte courant (quasi-noir)
\definecolor{colSlate}{RGB}{83, 98, 118}       % gris ardoise (méta, en-têtes)
\definecolor{colRule}{RGB}{203, 213, 225}      % filets clairs
\definecolor{colJuryBg}{RGB}{255, 248, 225}    % fond crème jury
\definecolor{colPresence}{RGB}{83, 98, 118}    % bloc 5 (présence dans les sujets)
\AtBeginDocument{\color{colInk}}

% ── Style encadré exercice + corrigé (banque d'exercices) ─────────────
\tcbset{
  exoBox/.style={
    breakable,
    enhanced,
    colback=colExo!8!white,
    colframe=colExo,
    colbacktitle=colExo,
    coltitle=white,
    fonttitle=\bfseries\small,
    left=4pt, right=4pt, top=3pt, bottom=3pt,
    before skip=5pt, after skip=5pt,
  }
}

% ── Bascule enonce / corrige (banque d'exercices) ─────────────────────
% \CAPETCORRIGE est defini en ligne de commande (pdflatex "\def\CAPETCORRIGE{1}\input{...}")
% pour compiler la variante corriges d'un meme fichier source ; par defaut (non defini),
% la compilation produit la variante enonces. \CAPETSHOW{enonce}{corrige} choisit la branche.
%
% TROIS MODES (refonte 2026-09-20, manuel de revision) :
%   - aucun drapeau   -> imprime l'enonce            (variante enonces, INCHANGE)
%   - \CAPETCORRIGE   -> imprime le corrige          (variante corriges, INCHANGE)
%   - \CAPETMANUEL    -> imprime l'enonce ET met le corrige de cote pour l'annexe
%
% Le mode manuel stocke chaque corrige dans une macro (\csgdef), sans ecriture
% fichier ni re-tokenisation : le contenu (circuitikz, accents, %) traverse intact.
% L'annexe les rejoue ensuite dans l'ordre via \CAPETANNEXE.
\usepackage{etoolbox}

\newcounter{CAPETexo}          % numero continu d'exercice dans le manuel
\setcounter{CAPETexo}{0}
\newcounter{CAPETannexe}       % compteur de relecture en annexe

% Titre courant de la notion en cours (pose par \CAPETNOTION), sert d'intertitre annexe
\newcommand{\CAPETnotioncode}{}
\newcommand{\CAPETnotiontitre}{}

% \CAPETNOTION{code}{titre} : ouvre une notion dans le manuel (memorise le libelle,
% et marque le rang du premier exercice de la notion pour l'intertitre d'annexe).
\newcommand{\CAPETNOTION}[2]{%
  \renewcommand{\CAPETnotioncode}{#1}%
  \renewcommand{\CAPETnotiontitre}{#2}%
  \csgdef{CAPETnotionstart\the\numexpr\value{CAPETexo}+1\relax}{#1 -- #2}%
}

% ── \degree robuste ───────────────────────────────────────────────────
% Certains fichiers d'exercices utilisent \degree. Ni le preambule ni siunitx
% ne le definissent, et \providecommand local ne suffit pas : sous hyperref,
% \degree existe deja comme commande factice (\ERROR), donc \providecommand
% s'abstient et l'appel echoue a la compilation du manuel (erreur invisible
% en compilation isolee). On le (re)definit donc inconditionnellement ici,
% apres le chargement des packages, pour tout le corpus.
\AtBeginDocument{\providecommand{\degree}{}\renewcommand{\degree}{\ensuremath{^\circ}}}

% Libelle decoratif : les fichiers d'exercices utilisent aussi \CAPETSHOW dans
% leur \fancyhead et leur bloc de titre, pour afficher "(enonces)" ou "(corriges)".
% Ces appels-la ne sont PAS des exercices : en mode manuel ils ne doivent ni
% consommer un numero, ni alimenter l'annexe. On les reconnait a leur premier
% argument, strictement egal au libelle ci-dessous.
\newcommand{\CAPETLIBELLEENONCES}{(\'enonc\'es)}

\ifdefined\CAPETMANUEL
  % --- Mode manuel : enonce en place, corrige differe vers l'annexe ---
  \newcommand{\CAPETSHOW}[2]{%
    \def\CAPETargun{#1}%
    \ifx\CAPETargun\CAPETLIBELLEENONCES
      #1%            appel decoratif : rendu tel quel, hors numerotation
    \else
      \CAPETSHOWEXO{#1}{#2}%
    \fi
  }
  \newcommand{\CAPETSHOWEXO}[2]{%
    \stepcounter{CAPETexo}%
    \csgdef{CAPETcorrige\theCAPETexo}{#2}%
    \edef\CAPETtmpnum{\theCAPETexo}%
    #1%
    \par\vspace{2pt}%
    {\footnotesize\itshape\color{colExo}%
     \hyperref[capet:cor:\CAPETtmpnum]{Corrige de l'exercice \CAPETtmpnum\ $\rightarrow$ p.~\pageref{capet:cor:\CAPETtmpnum}}%
     \label{capet:eno:\CAPETtmpnum}}%
  }
  % \CAPETANNEXE : rejoue tous les corriges captures, dans l'ordre, avec intertitres
  \newcommand{\CAPETANNEXE}{%
    \setcounter{CAPETannexe}{0}%
    \loop\ifnum\value{CAPETannexe}<\value{CAPETexo}%
      \stepcounter{CAPETannexe}%
      \edef\CAPETcurnum{\theCAPETannexe}%
      \ifcsname CAPETnotionstart\CAPETcurnum\endcsname
        \par\medskip
        {\normalsize\bfseries\color{colExo}\csuse{CAPETnotionstart\CAPETcurnum}}%
        \par\nobreak\smallskip
      \fi
      \begin{tcolorbox}[exoBox, breakable,
        title={Corrige de l'exercice \CAPETcurnum\
               \hfill\normalfont\footnotesize
               \hyperref[capet:eno:\CAPETcurnum]{enonce p.~\pageref{capet:eno:\CAPETcurnum}}}]%
      \label{capet:cor:\CAPETcurnum}%
      \csuse{CAPETcorrige\CAPETcurnum}%
      \end{tcolorbox}%
    \repeat
  }
\else
  % --- Modes historiques, comportement strictement inchange ---
  \newcommand{\CAPETSHOW}[2]{\ifdefined\CAPETCORRIGE #2\else #1\fi}
  \newcommand{\CAPETANNEXE}{}
\fi

% ── Difficulté 1 à 3 (étoiles pleines, amssymb) ───────────────────────
\newcommand{\difficulte}[1]{%
  \ifcase#1\relax
  \or $\bigstar$\;niveau 1/3%
  \or $\bigstar\bigstar$\;niveau 2/3%
  \or $\bigstar\bigstar\bigstar$\;niveau 3/3%
  \fi
}

% ── Style encadré extrait de document technique (banque d'exercices) ──
\tcbset{
  docTech/.style={
    breakable,
    enhanced,
    colback=colGrayLight,
    colframe=colSeparator,
    fonttitle=\bfseries\small\itshape,
    coltitle=black,
    colbacktitle=colGrayLight,
    left=4pt, right=4pt, top=3pt, bottom=3pt,
    before skip=4pt, after skip=4pt,
  }
}

% ── Style tcolorbox de section principale ────────────────────────────
% Refonte 2026 : bandeau titre coloré arrondi + fond très pâle, sans cadre.
\tcbset{
  ficheSection/.style={
    breakable, enhanced jigsaw,
    frame hidden, boxrule=0pt, arc=1.8mm,
    colback=#1!4!white,
    colbacktitle=#1,
    coltitle=white,
    fonttitle=\sffamily\bfseries,
    left=7pt, right=7pt, top=5pt, bottom=5pt,
    toptitle=3pt, bottomtitle=3pt,
    before skip=10pt, after skip=8pt,
  }
}

% ── Style encadré question type ──────────────────────────────────────
% Carte blanche, bandeau sombre, liseré très clair.
\tcbset{
  questiontype/.style={
    breakable, enhanced jigsaw,
    colback=white,
    frame hidden, boxrule=0pt, arc=1.8mm,
    colbacktitle=colElec!78!black,
    coltitle=white,
    fonttitle=\sffamily\bfseries\small,
    left=7pt, right=7pt, top=4pt, bottom=4pt,
    toptitle=2.5pt, bottomtitle=2.5pt,
    before skip=8pt, after skip=8pt,
    borderline={0.6pt}{0pt}{colElec!12!white},
  }
}

% ── Style encadré jury ───────────────────────────────────────────────
% Bandeau ambre, fond crème, sans cadre.
\tcbset{
  jury/.style={
    breakable, enhanced jigsaw,
    colback=colJuryBg,
    frame hidden, boxrule=0pt, arc=1.8mm,
    colbacktitle=colJuryFrame,
    coltitle=white,
    fonttitle=\sffamily\bfseries,
    left=7pt, right=7pt, top=4pt, bottom=4pt,
    toptitle=3pt, bottomtitle=3pt,
    before skip=10pt, after skip=8pt,
  }
}

% ── En-tête / pied de page communs (refonte 2026) ─────────────────────
% #1 code fiche, #2 titre court, #3 couleur domaine
\newcommand{\ficheHeaders}[3]{%
  \fancyhead[L]{\small\sffamily{\bfseries\textcolor{#3}{#1}}\hspace{0.7em}\textcolor{colSlate}{#2}}%
  \fancyhead[R]{\small\sffamily\textcolor{colSlate}{CAPET SII IE\, --- \,Session 2026}}%
  \fancyfoot[C]{\small\sffamily\textcolor{colSlate}{\thepage\,/\,\pageref{LastPage}}}%
  \fancyhead[C]{}\fancyfoot[L]{}\fancyfoot[R]{}%
}

% ── Bloc en-tête de fiche (BLOC 0, refonte 2026) ──────────────────────
% #1 code  #2 titre  #3 couleur domaine  #4 domaine  #5 ligne présence  #6 sous-thèmes
\newcommand{\ficheEntete}[6]{%
  \begin{tcolorbox}[enhanced, breakable,
    frame hidden, boxrule=0pt, arc=2.5mm,
    interior style={left color=#3, right color=#3!75!black},
    left=9pt, right=9pt, top=7pt, bottom=7pt,
    before skip=0pt, after skip=6pt]
    {\color{white}\sffamily\bfseries\LARGE #1\quad #2}\\[4pt]
    {\color{#3!22!white}\sffamily\small \MakeUppercase{#4}\quad\textbar\quad #5}\\[1pt]
    {\color{#3!35!white}\sffamily\footnotesize Sous-thèmes : #6}
  \end{tcolorbox}%
}

% ── Mise en page ─────────────────────────────────────────────────────
\pagestyle{fancy}
\fancyhf{}
\renewcommand{\headrulewidth}{0.5pt}
\renewcommand{\footrulewidth}{0pt}
\AtBeginDocument{%
  \renewcommand{\headrule}{{\color{colRule}\hrule width\headwidth height 0.5pt}}%
}

\setlist[itemize]{leftmargin=1.2em, itemsep=1pt, topsep=2pt, parsep=0pt}
\setlist[enumerate]{leftmargin=1.2em, itemsep=3pt, topsep=2pt, parsep=0pt}

\setlength{\parskip}{2pt}
\setlength{\columnsep}{10pt}

\def\CAPETBEGINFILE{\begin{document}}
\def\CAPETENDFILE{\end{document}}
\else
\def\CAPETBEGINFILE{}
\def\CAPETENDFILE{}
\fi

\ficheHeaders{T06}{Transformateur}{colElec}

\CAPETBEGINFILE

% -----------------------------------------------------------------------
%  BLOC EN-TETE
% -----------------------------------------------------------------------

\ficheEntete{T06}{Transformateur (monophasé et triphasé)}{colElec}%
  {électrotechnique}%
  {\textbf{Présence :} \textbf{13/13} sujets (2013-2026) - thème \emph{omniprésent}}%
  {rapport de transformation - essai à vide ($R_f$, $X_m$) - essai court-circuit ($R_s$, $\varepsilon_{cc}$) - couplages Yy/Yd - chute de tension - rendement}

% -----------------------------------------------------------------------
%  1. RAPPELS DE COURS
% -----------------------------------------------------------------------

\begin{tcolorbox}[
  ficheSection=colElec,
  title={1. Rappels de cours},
]

% -- 1.1 Rapport de transformation ------------------------------------
\textbf{1.1 Rapport de transformation}

\smallskip
\begin{center}
\renewcommand{\arraystretch}{1.45}
\begin{tabular}{lll}
\toprule
\textbf{Grandeur} & \textbf{Expression} & \textbf{Remarque} \\
\midrule
Rapport de transformation & $m = \dfrac{N_2}{N_1} = \dfrac{V_2}{V_1} = \dfrac{I_1}{I_2}$ & Monophasé (tensions à vide) \\[8pt]
Rapport (triphasé) & $m = \dfrac{U_{2n}}{U_{1n}}$ & Tensions composées étoile-étoile \\[8pt]
Tension secondaire à vide & $U_s = m \cdot U_1$ & Sans charge \\[4pt]
Couplage Yy & $U_{composee} = \sqrt{3}\cdot U_{simple}$ & Rapport $m$ inchangé \\[4pt]
Couplage Yd & Décalage de $30^\circ$ prim./sec. & Élimine harmoniques rang 3 \\
\bottomrule
\end{tabular}
\end{center}

\medskip
\hrule\vspace{4pt}

% -- 1.2 Schema equivalent - essai a vide -----------------------------
\textbf{1.2 Essai à vide - pertes fer}

\smallskip
À vide ($I_2 = 0$) : le schéma équivalent se réduit à $R_f$ et $X_m$ en parallèle, alimentés par $V_1$.
\[
  P_0 = \frac{V_1^2}{R_f} \quad\Rightarrow\quad R_f = \frac{V_1^2}{P_0}
  \qquad;\qquad
  Q_0 = \frac{V_1^2}{X_m} \quad\Rightarrow\quad X_m = \frac{V_1^2}{Q_0}
\]
\[
  S_0 = V_1 \cdot I_0 \qquad;\qquad Q_0 = \sqrt{S_0^2 - P_0^2}
\]

\smallskip
\textbf{Schéma équivalent de l'essai à vide :}
\begin{center}
\begin{circuitikz}[european, scale=0.9]
  \draw (0,0) to[V, l=$V_1$] (0,2.6) node[circ]{}
        to[short] (2,2.6) node[circ]{}
        to[R, l=$R_f$] (2,0) node[circ]{}
        -- (0,0) node[circ]{};
  \draw (2,2.6) to[short] (3.4,2.6) node[circ]{}
        to[cute inductor, l=$X_m$] (3.4,0) node[circ]{}
        -- (2,0);
\end{circuitikz}
\end{center}
$R_f$ modélise les pertes fer (dissipatives), $X_m$ le courant magnétisant (réactif). Les deux branches sont en parallèle sous $V_1$, d'où $S_0^2 = P_0^2 + Q_0^2$.

\textit{$P_0$ = pertes fer (mesurées au wattmètre côté primaire). $Q_0$ est toujours positif (charge inductive).}

\medskip
\hrule\vspace{4pt}

% -- 1.3 Essai en court-circuit - pertes cuivre -----------------------
\textbf{1.3 Essai en court-circuit - pertes cuivre}

\smallskip
\[
  P_{cc} = R_s \cdot I_{cc}^2 \quad\Rightarrow\quad R_s = R_1 + m^2 R_2 \quad \text{(résistances ramenées au primaire)}
\]
\[
  \varepsilon_{cc} = \frac{V_{cc}}{V_{1n}} \quad (\text{valeur typique : } 4\,\% \text{ à } 10\,\%)
\]
Chute de tension relative en charge :
\[
  \Delta U \approx \varepsilon_R \cos\varphi + \varepsilon_X \sin\varphi \qquad \text{avec}\quad
  \varepsilon_R = \frac{P_{cc}}{S_n},\quad \varepsilon_X \approx \varepsilon_{cc}
\]

\medskip
\hrule\vspace{4pt}

% -- 1.4 Rendement ----------------------------------------------------
\textbf{1.4 Rendement}

\smallskip
\[
  \eta = \frac{P_2}{P_1} = \frac{P_2}{P_2 + P_{fer} + P_{cuivre}}
\]
Maximum lorsque $P_{fer} = P_{cuivre}$ ; en pratique, le taux de charge optimal dépend des pertes nominales du transformateur et doit être calculé cas par cas.

\medskip\hrule\vspace{4pt}

% Complément du référentiel - 2026-11-10 (sans TikZ : les ids de schémas restent stables)
\textbf{1.5 Circuit magnétique et électromagnétisme (complément)}

\smallskip
Théorème d'Ampère sur la ligne de champ moyenne : $N I = H\,\ell$. Loi d'Hopkinson : $N I = \mathcal{R}\,\Phi$ avec $\mathcal{R} = \ell/(\mu_0\mu_r S)$ ; l'entrefer, de perméabilité $\mu_0$, porte souvent l'essentiel de la réluctance. Loi de Faraday : $e = -N\,\mathrm{d}\Phi/\mathrm{d}t$. En régime sinusoïdal, la formule de Boucherot $V_{eff} = 4{,}44\,N\,f\,B_{max}\,S$ fixe le nombre de spires pour que l'induction crête reste sous la saturation (1{,}2 à 1{,}5\,T dans les tôles usuelles).

\end{tcolorbox}

% -----------------------------------------------------------------------
%  2. FORMULAIRE SYNTHETIQUE
% -----------------------------------------------------------------------

\begin{tcolorbox}[
  ficheSection=colElec,
  title={2. Formulaire synthétique},
]
\small
\begin{multicols}{2}
\renewcommand{\arraystretch}{1.45}
\begin{tabular}{@{}p{4.5cm}ll@{}}
\toprule
\textbf{Formule} & \textbf{Usage} & \textbf{U.} \\
\midrule
\multicolumn{3}{@{}l}{\footnotesize\textit{Rapport de transformation}} \\
$m = N_2/N_1 = U_2/U_1 = I_1/I_2$ & Rapport & - \\
$U_s = m \cdot U_1$ & Tension sec. & V \\
\midrule
\multicolumn{3}{@{}l}{\footnotesize\textit{Essai à vide}} \\
$S_0 = V_1 I_0$ & Puiss. app. & VA \\
$Q_0 = \sqrt{S_0^2 - P_0^2}$ & Puiss. réact. & var \\
$R_f = V_1^2 / P_0$ & Rés. perte fer & $\Omega$ \\
$X_m = V_1^2 / Q_0$ & React. mag. & $\Omega$ \\
\bottomrule
\end{tabular}

\columnbreak

\renewcommand{\arraystretch}{1.45}
\begin{tabular}{@{}p{4.5cm}ll@{}}
\toprule
\textbf{Formule} & \textbf{Usage} & \textbf{U.} \\
\midrule
\multicolumn{3}{@{}l}{\footnotesize\textit{Essai court-circuit}} \\
$R_s = P_{cc}/I_{cc}^2$ & Rés. ramenée & $\Omega$ \\
$\varepsilon_{cc} = V_{cc}/V_{1n}$ & Tension rel. & - \\
$\Delta U \approx \varepsilon_R\cos\varphi + \varepsilon_X\sin\varphi$ & Chute tens. & - \\
\midrule
\multicolumn{3}{@{}l}{\footnotesize\textit{Rendement}} \\
$\eta = P_2/(P_2 + P_{fer} + P_{cu})$ & Rendement & - \\
$P_{fer} = P_{cu}$ & Max. $\eta$ & - \\
$\eta_{typ} > 95\,\%$ & Valeur usuelle & - \\
\bottomrule
\end{tabular}
\end{multicols}
\normalsize
\end{tcolorbox}

% -----------------------------------------------------------------------
%  3. QUESTIONS TYPES
% -----------------------------------------------------------------------

\begin{tcolorbox}[
  ficheSection=colElec,
  title={3. Questions types - extraites des sujets réels},
]

% -- QT1 -------------------------------------------------------------

\begin{tcolorbox}[questiontype,
  title={QT1 - Rapport de transformation - redresseur P3 vers 600\,V DC \hfill\normalfont\footnotesize E1 sujet 2026}
]
\textbf{Énoncé.}
Le réseau triphasé fournit $230/400\,\mathrm{V}$. Un transformateur étoile-étoile alimente un redresseur P3 (3 diodes, cathodes communes). La tension moyenne de sortie est $\bar{u}_s = \dfrac{3\sqrt{6}}{2\pi} \cdot V_{eff,simple}$. Le cahier des charges impose $\bar{u}_s = 600\,\mathrm{V}$. Calculer $V_{s,simple}$ puis le rapport de transformation $m$.

\tcbline

\textbf{Corrigé.}

\textbf{Étape 1 — Tension simple secondaire nécessaire.}
\[
  V_{s,simple} = \frac{600 \times 2\pi}{3\sqrt{6}} = \frac{600 \times 6{,}283}{7{,}348} \approx 513\,\mathrm{V}
\]

\textbf{Étape 2 — Rapport de transformation.}
\[
  m = \frac{V_{s,simple}}{V_{1,simple}} = \frac{513}{230} \approx \mathbf{2{,}23}
\]

\textit{Rappel : en couplage étoile-étoile, le rapport $m$ s'applique indifféremment aux tensions simples ou composées.}

\textbf{Conclusion.} Le rapport $m \approx 2{,}23$ permet d'atteindre exactement les $600\,\mathrm{V}$ imposés par le cahier des charges pour l'option de redressement P3.
\end{tcolorbox}

\vspace{4pt}

% -- QT2 -------------------------------------------------------------

\begin{tcolorbox}[questiontype,
  title={QT2 - Rapport de transformation - redresseur PD3 vers 600\,V DC \hfill\normalfont\footnotesize E1 sujet 2026}
]
\textbf{Énoncé.}
Pour la solution PD3 (pont parallèle double triphasé, 6 diodes), la tension moyenne est $\bar{u}_s = 0{,}95 \cdot U_M$ où $U_M$ est l'amplitude des tensions composées secondaires. Pour obtenir $\bar{u}_s = 600\,\mathrm{V}$ avec un primaire $400\,\mathrm{V}$ composé, calculer $U_M$, puis $U_{s,composee}$, puis $m$.

\tcbline

\textbf{Corrigé.}

\textbf{Étape 1 — Amplitude des tensions composées.}
\[
  U_M = \frac{600}{0{,}95} \approx 631\,\mathrm{V}
\]

\textbf{Étape 2 — Valeur efficace secondaire.}
\[
  U_{s,composee} = \frac{U_M}{\sqrt{2}} = \frac{631}{1{,}414} \approx 446\,\mathrm{V}
\]

\textbf{Étape 3 — Rapport de transformation.}
\[
  m = \frac{U_{s,composee}}{U_{1,composee}} = \frac{446}{400} \approx \mathbf{1{,}12}
\]

\textit{Comparaison : P3 nécessite $m \approx 2{,}23$ contre $1{,}12$ pour PD3 - le PD3 exploite mieux le transformateur.}

\textbf{Conclusion.} La solution PD3 (dimensionnée avec $m \approx 1{,}12$) atteint le même objectif de $600\,\mathrm{V}$ DC avec un rapport de transformation près de deux fois plus faible que le P3 : elle est préférable si le composant transformateur doit être minimisé.
\end{tcolorbox}

\vspace{4pt}

% -- QT3 -------------------------------------------------------------

\begin{tcolorbox}[questiontype,
  title={QT3 - Essai à vide - calcul de $R_f$ et $X_m$ \hfill\normalfont\footnotesize E2 BAC3 2026, Q.~A12--A13 (MAS à $g\approx0$)}
]
\textbf{Énoncé.}
Lors de l'essai à vide d'une machine asynchrone triphasée ($g \approx 0$), le schéma équivalent d'une phase se réduit à $R_f$ et $X_m$ en parallèle. Données du sujet : tension simple nominale $V_n = \SI{230}{\volt}$, courant de ligne à vide $I_0 = \SI{42.5}{\ampere}$, puissance active à vide $P_0 = \SI{7.82}{\kilo\watt}$ (triphasée). Calculer $S_0$, puis en déduire $R_f$ et $X_m$.

\tcbline

\textbf{Corrigé.} \textit{[corrigé reconstitué]}

\textbf{Étape 1 — Puissance apparente à vide (triphasée).}
\[
  S_0 = 3\,V_n I_0 = 3 \times 230 \times 42{,}5 = 29\,325\,\mathrm{VA} \approx \SI{29.3}{\kilo\volt\ampere}
\]

\textbf{Étape 2 — Résistance des pertes fer.}
$P_0$ étant la puissance triphasée, chaque phase dissipe $P_0/3$ dans $R_f$ :
\[
  R_f = \frac{V_n^2}{P_0/3} = \frac{3 \times 230^2}{7\,820} = \frac{158\,700}{7\,820} \approx \mathbf{20{,}3\,\Omega}
\]

\textbf{Étape 3 — Puissance réactive et réactance magnétisante.}
\[
  Q_0 = \sqrt{S_0^2 - P_0^2} = \sqrt{29\,325^2 - 7\,820^2} \approx 28\,263\,\mathrm{var}
  \qquad
  X_m = \frac{3\,V_n^2}{Q_0} = \frac{158\,700}{28\,263} \approx \mathbf{5{,}6\,\Omega}
\]

\textbf{Conclusion.} On vérifie $S_0 > P_0$ (cohérence physique) et $\cos\varphi_0 = P_0/S_0 \approx 0{,}27$, valeur faible caractéristique d'un essai à vide dominé par le courant magnétisant : le modèle parallèle $R_f$--$X_m$ est validé et directement exploitable pour le bilan de puissance réactive de l'installation.
\end{tcolorbox}

\end{tcolorbox}

% -----------------------------------------------------------------------
%  4. ATTENTES DU JURY
% -----------------------------------------------------------------------

\begin{tcolorbox}[
  jury,
  title={4. Attentes du jury - Erreurs récurrentes et points de vigilance},
]

\begin{enumerate}[label=\textbf{\arabic*.}, leftmargin=*, itemsep=5pt]

  \item \textbf{Confondre tensions simples et composées.}
  En couplage étoile-étoile : $m = U_{s,comp}/U_{1,comp} = V_{s,simple}/V_{1,simple}$ - les deux sont équivalents.
  En couplage étoile-triangle, $m$ relie la tension simple primaire à la tension composée secondaire (ou vice-versa).

  \item \textbf{Oublier le facteur $\sqrt{3}$.}
  $U_{composee} = \sqrt{3} \cdot V_{simple}$ en étoile. Ne pas confondre les deux niveaux de tension.

  \item \textbf{Inverser essai à vide et essai court-circuit.}
  À vide : pertes \textbf{fer} ($R_f$, $X_m$), courant $I_0$ faible, tension nominale.
  Court-circuit : pertes \textbf{cuivre} ($R_s$), tension réduite $V_{cc}$, courant nominal.

  \item \textbf{Signe de $Q_0$.}
  $Q_0 = \sqrt{S_0^2 - P_0^2}$ est toujours \textbf{positif} (charge inductive).
  L'erreur classique est d'écrire $Q_0 = \sqrt{P_0^2 - S_0^2}$, ce qui donne un résultat imaginaire.

  \item \textbf{Rapport $m$ et redresseur.}
  Le transformateur en amont d'un redresseur a un $m$ déterminé par la tension DC voulue.
  Calculer d'abord $U_{s,AC}$ nécessaire depuis la formule $\bar{U}_{DC}$, puis $m = U_s/U_1$.

  \item \textbf{Justifier la démarche avant le calcul.}
  Énoncer le modèle, les hypothèses (fréquence fixe, circuit magnétique linéaire), puis calculer.

\end{enumerate}

\medskip
\begin{tcolorbox}[
  colback=colJuryFrame!12!colJury,
  colframe=colJuryFrame,
  left=4pt, right=4pt, top=2pt, bottom=2pt,
  fontupper=\small\itshape
]
Rapport jury 2025 : «~De nombreux candidats ne maitrisent pas les lois basiques de l'électricité (la loi des nœuds, la loi des mailles, les calculs de puissance ou encore la loi d'Ohm).~»\\
Rapport jury 2025 : «~Le jury apprécie les copies des candidats qui justifient ou expliquent les démarches adoptées pour répondre aux questions posées.~»\\
Point de vigilance récurrent (non cité verbatim par les rapports) : la confusion entre couplages étoile/triangle et rapports de transformation associés — vérifier si $m$ porte sur les tensions simples ou composées selon le couplage.
\end{tcolorbox}

\end{tcolorbox}

% -----------------------------------------------------------------------
%  PIED - PRESENCE DANS LES SUJETS
% -----------------------------------------------------------------------

\begin{tcolorbox}[
  ficheSection=colPresence,
  title={Présence dans les sujets analysés (2013-2026)},
]
\begin{center}
\footnotesize\renewcommand{\arraystretch}{1.15}
\begin{tabular}{lccccccccccccc}
\toprule
\textbf{Année} & 13 & 14 & 15 & 16 & 17 & 18 & 19 & 20 & 21 & 22 & 23 & 24 & 25 \\
\midrule
E1 (commune)   & $\checkmark$ & $\checkmark$ & - & $\checkmark$ & - & $\checkmark$ & - & - & $\checkmark$ & - & $\checkmark$ & - & $\checkmark$ \\
E2 discipl. IE & - & $\checkmark$ & $\checkmark$ & - & $\checkmark$ & - & $\checkmark$ & $\checkmark$ & - & $\checkmark$ & - & $\checkmark$ & - \\
\bottomrule
\end{tabular}
\end{center}
\smallskip
\footnotesize\textbf{Sous-thèmes fréquents :}
rapport de transformation (10/13) - essai à vide $R_f$/$X_m$ (6/13) - couplages triphasés Yy/Yd (5/13) - redresseur + transformateur amont (4/13) - rendement (3/13)

\medskip
\textbf{Ordres de grandeur (repères rapides).}
\begin{itemize}[leftmargin=1.3em]
  \item Réseau BT FR : tension phase-neutre typique \SI{230}{\volt}, tension phase-phase typique \SI{400}{\volt}, fréquence \SI{50}{\hertz}.
  \item Courants d'étude fréquents en sujet : de quelques \si{\ampere} à quelques dizaines de \si{\ampere}.
  \item En conversion d'énergie, un rendement attendu reste généralement entre 0{,}8 et 0{,}98 selon la technologie.
\end{itemize}

\medskip
\textbf{Synthèse express (à refaire sans support).}
\begin{enumerate}[leftmargin=1.4em]
  \item Poser les hypothèses du modèle et les conventions de signe.
  \item Écrire 2 à 4 relations clés avec unités.
  \item Réaliser un calcul numérique complet et vérifier l'ordre de grandeur.
  \item Justifier le résultat par un argument technique de niveau CAPET.
\end{enumerate}

\end{tcolorbox}

\CAPETENDFILE
